% Append-only section fragment. Paste below the preceding section in the master document.
% This fragment intentionally contains no preamble and no document boundaries.

\section{Evidence-Backed Programme Cast Companion}
\label{sec:programme-cast-companion}

Pilot Fabric 0.42 begins the next Live Companion expansion by answering programme
cast and character questions from verified public evidence. The capability is
designed around a strict distinction: knowing the cast of a verified programme is
not the same as identifying a person visible in the current scene.

\subsection{User Capability}

The initial conversational interface supports three request classes:

\begin{description}
  \item[Programme cast] Requests such as ``Pilot, who is in this?'' and
        ``Pilot, show me the cast'' return a bounded list of main cast members and
        their characters.
  \item[Exact character] A request such as ``Pilot, who plays Queen Elizabeth II?''
        returns an actor only when the character name has one exact normalized match
        in the verified cast evidence.
  \item[Visible actor] A request such as ``Pilot, who is that actor?'' receives the
        explicit statement that Pilot cannot identify the person currently visible.
        Pilot then offers the verified main cast as a useful, non-biometric alternative.
\end{description}

The short answer is spoken without leaving the programme. The private phone shows
up to eight cast entries, the programme title, attribution and the safety boundary.

\subsection{Evidence Pipeline}

The runtime uses the following deterministic sequence:

\begin{enumerate}
  \item The voice parser maps the bounded request to the
        \texttt{programme\_cast} intent.
  \item The television gateway observes permitted device state without changing
        application, input or playback.
  \item The multi-frame perception timeline supplies the current stable programme
        title and its identity provenance.
  \item If the existing provider metadata contains one exact TVmaze catalogue
        identifier, that identifier is reused.
  \item Otherwise, Pilot performs a bounded title search and accepts an identifier
        only when exactly one result has the same normalized title.
  \item Pilot requests the TVmaze main-cast endpoint for that exact programme.
  \item Person and character names are normalized into a bounded evidence record.
  \item The answer is produced deterministically from that record; no language
        model is required to recall or invent cast information.
\end{enumerate}

The TVmaze main-cast endpoint orders results by the importance derived from the
number of character appearances. Pilot treats this as a programme-level cast list,
not proof that any member appears in the current episode, frame or scene.

\subsection{No-Face-Guessing Boundary}

This capability does not run face detection, facial recognition, face embeddings,
biometric matching, celebrity recognition or visual person search. It does not
compare the television image with public photographs. The response record always
contains explicit false values for \texttt{face\_recognition\_used} and
\texttt{visible\_person\_identified}.

Consequently, ``who is that actor?'' is interpreted as a request for assistance,
not authority to identify the visible individual. Pilot says that it cannot identify
the person and lists the verified cast. A future product may let the owner select a
character name or inspect permitted episode credits, but it must preserve the same
prohibition on ungrounded face identity.

\subsection{Character Matching}

Character names are Unicode-normalized, stripped of punctuation and compared in a
case-insensitive exact form. Substring or fuzzy similarity is not sufficient. If no
exact character exists, Pilot states that it could not verify the requested role.
This prevents a similarly named character from being silently mapped to the wrong
actor.

\subsection{Caching, Limits and Attribution}

Public cast responses are limited to 500 kilobytes and a ten-second request timeout.
Pilot processes at most forty upstream cast rows, retains at most twenty normalized
entries and projects no more than twelve into the evidence record. The private phone
shows no more than eight.

Cast evidence is cached on the mother brain for twenty-four hours. The cache is
limited to fifty programme identifiers and uses owner-only file permissions. Oldest
entries are discarded when the bound is exceeded. Each record retains the source
URL and the attribution ``Cast data: TVmaze (CC BY-SA)''. Production redistribution
must continue to respect TVmaze attribution and share-alike requirements.

\subsection{Privacy and Playback Preservation}

The public request contains only the exact catalogue identifier or a stable text
title needed to establish that identifier. Pilot sends no source pixels, face data,
audio, subtitles, provider credentials, household account identifiers, viewing
history or private persona data. Protected programme media is neither copied nor
replaced. The television remains on its existing surface throughout the operation.

The audit receipt records the lookup identifier, programme title, catalogue
identifier, question class, identity route and evidence hash. It explicitly records
that face recognition was not used, visible-person identity was not established and
playback was preserved.

\subsection{Failure Behaviour}

Pilot fails closed when no programme title exists, the title is not stable, the
catalogue search is ambiguous, the catalogue cannot be reached, the response exceeds
its size bound, the cast record is empty or an exact character is absent. A failed
lookup never falls through to model memory, web-image matching or a guessed answer.

\subsection{Verification}

Automated tests cover verified cast retrieval, reuse of an existing exact catalogue
identifier, twenty-four-hour cache reuse, exact character matching, visible-actor
refusal language, the absence of face-recognition claims, playback preservation,
failure without a stable programme identity and deterministic voice routing. Pilot
Fabric 0.42 passes 314 automated Fabric tests.

\subsection{Remaining Live Companion Work}

This closes the first cast-lookup capability but not the complete Live Companion
section. Remaining work includes broader person credits, episode-specific guest
cast, evidence-backed filming locations and original sources, culturally specific
references, child-friendly reading levels, accessibility summaries, permitted
subtitle projection, richer authenticated live-event adapters and native compact
overlays where a television platform grants the required authority.
